%================================================================
\chapter{Assessment and Progress Record}
\label{ch:assessment}
%================================================================

\begin{Form}

This chapter is a fillable PDF form. Each section corresponds to one
category of prerequisite: Required, Essential, Desirable, and Optional.
Checkboxes record binary completion; text fields record written
responses. Use this chapter to track your progress through the program.

\vs
\noindent
\textbf{Saving your responses requires Adobe Acrobat Reader} (free for
Windows, macOS, and Linux; available at
\url{https://get.adobe.com/reader/}). Other viewers behave as follows:

\begin{itemize}
    \item \textbf{macOS Preview:} displays form fields and allows you
          to type in them, but \emph{does not save the entered data}.
          If you close and reopen the file in Preview, your responses
          will be lost.
    \item \textbf{Google Chrome (built-in PDF viewer):} displays and
          accepts input in form fields, but \emph{does not save the
          data}. Printing to PDF from Chrome will not preserve field
          contents either.
    \item \textbf{Most Linux PDF viewers} (Evince, Okular, Zathura):
          \emph{ignore form fields entirely}. Fields will not appear
          interactive.
    \item \textbf{Adobe Acrobat Reader:} fully saves, prints, and
          exports form data. Use \textbf{File $>$ Save} (not Save As)
          after completing each section.
\end{itemize}

\vs
\noindent
\textbf{How to interpret the four categories:}

\begin{description}
    \item[Required] Items directly related to the Data Challenge.
          These must be completed to participate in the program.
    \item[Essential] Environment setup items: Python, Visual Studio
          Code, Git, the required libraries, and a successful API
          call to a commercial LLM. These must be completed before
          the first session.
    \item[Desirable] Technical understanding of tools covered in the
          prerequisite chapters (CLI, programming, virtual
          environments, Git, \LaTeX{}). Completing these demonstrates
          readiness to work independently.
    \item[Optional] LLM infrastructure setup (Ollama). Not required
          for the data challenge; relevant for students who want to
          run local models.
\end{description}
\vs

%================================================================
\paragraph{Required: Data Challenge}
\label{sec:assess-required}
%================================================================

The following items require engaging directly with the NCHS vital
statistics dataset. Items R3 and R4 require running Python code.
See Chapter~\ref{ch:analysis-activities} for the analysis context.

\vs

\noindent\textbf{R1.} Why is the date and time of birth no longer
recorded in NCHS data after 1988? What does this decision tell us
about how data governance priorities can change over time?
(See Chapter~\ref{ch:analysis-activities}.)\\[4pt]
\TextField[name=r1,width=\linewidth,height=3cm,multiline=true,
           bordercolor=black]{}

\vs

\noindent\textbf{R2.} Why was the flat file format used for NCHS
birth records? What are the tradeoffs of this format compared to a
relational database or a structured format like CSV with a header?
(See Chapter~\ref{ch:analysis-activities}.)\\[4pt]
\TextField[name=r2,width=\linewidth,height=3cm,multiline=true,
           bordercolor=black]{}

\vs

\noindent\textbf{R3.} How many live births occurred in Texas in 1969
from mothers residing in Texas? State your count and describe the
filtering logic you used.
(See Chapter~\ref{ch:analysis-activities}.)\\[4pt]
\TextField[name=r3,width=\linewidth,height=3cm,multiline=true,
           bordercolor=black]{}

\vs

\noindent\textbf{R4.} What socioeconomic variable in the NCHS dataset
shows the strongest relationship with birth weight in your analysis?
Briefly describe how you identified this relationship.
(See Chapter~\ref{ch:analysis-activities}.)\\[4pt]
\TextField[name=r4,width=\linewidth,height=3cm,multiline=true,
           bordercolor=black]{}

%================================================================
\paragraph{Essential: Environment and API Setup}
\label{sec:assess-essential}
%================================================================

Check each box when the corresponding verification step succeeds.
All six items should be checked before the first program session.

\vs

\noindent\CheckBox[name=e1,width=0.4cm,height=0.4cm]{}~\textbf{E1.}
\texttt{python --version} returns a version number in the terminal.
(See Chapter~\ref{ch:environment-setup}, Section~\ref{sec:python}.)

\vs

\noindent\CheckBox[name=e2,width=0.4cm,height=0.4cm]{}~\textbf{E2.}
\texttt{code .} opens Visual Studio Code from a terminal window.
(See Chapter~\ref{ch:environment-setup}, Section~\ref{sec:vscstart}.)

\vs

\noindent\CheckBox[name=e3,width=0.4cm,height=0.4cm]{}~\textbf{E3.}
\texttt{git --version} returns a version number.
(See Chapter~\ref{ch:environment-setup}, Section~\ref{sec:git}.)

\vs

\noindent\CheckBox[name=e4,width=0.4cm,height=0.4cm]{}~\textbf{E4.}
The five required libraries (\texttt{numpy}, \texttt{pandas},
\texttt{matplotlib}, \texttt{scipy}, \texttt{jupyter}) import
without errors in your virtual environment.
(See Chapter~\ref{ch:python-environments},
Section~\ref{sec:project-libraries}.)

\vs

\noindent\CheckBox[name=e5,width=0.4cm,height=0.4cm]{}~\textbf{E5.}
Running \texttt{agentClaude.py}, \texttt{agentGPT.py}, or
\texttt{agentGroq.py} from \texttt{resources/unit\_7/FOO/} returns
a response from a commercial LLM.
(See Chapter~\ref{ch:ollama-infrastructure},
Section~\ref{sec:apikeys}.)

\vs

\noindent\CheckBox[name=e6,width=0.4cm,height=0.4cm]{}~\textbf{E6.}
Running \texttt{pdflatex hello.tex} produces a \texttt{hello.pdf}
file.
(See Chapter~\ref{ch:latex}, Section~\ref{sec:first-latex-doc}.)

%================================================================
\paragraph{Desirable: Technical Skills}
\label{sec:assess-desirable}
%================================================================

These items demonstrate working understanding of the tools introduced
in the prerequisite chapters. Written responses need not be long;
two or three sentences each are sufficient.

\vs

\noindent\textbf{D1.} What is the current working directory when you
first open a terminal? Write the command you would use to navigate
to a folder named \texttt{DAIR3-Workshop} on your Desktop.
(See Chapter~\ref{ch:cli-scripting}, Section~\ref{sec:what-is-a-file}.)\\[4pt]
\TextField[name=d1,width=\linewidth,height=2.5cm,multiline=true,
           bordercolor=black]{}

\vs

\noindent\textbf{D2.} What happens when you type \texttt{hello.py}
in a terminal without \texttt{python} in front of it? Why does
this happen, and how do you fix it?
(See Chapter~\ref{ch:intro-prog}, Section~\ref{sec:plaintext-first}.)\\[4pt]
\TextField[name=d2,width=\linewidth,height=2.5cm,multiline=true,
           bordercolor=black]{}

\vs

\noindent\textbf{D3.} Why does adding 0.1 ten times in Python not
equal exactly 1.0? What underlying technical limitation causes this?
(See Chapter~\ref{ch:intro-prog}.)\\[4pt]
\TextField[name=d3,width=\linewidth,height=2.5cm,multiline=true,
           bordercolor=black]{}

\vs

\noindent\textbf{D4.} What is the difference between a Python virtual
environment and the global Python installation? Give one scenario
where you would need more than one virtual environment.
(See Chapter~\ref{ch:python-environments}.)\\[4pt]
\TextField[name=d4,width=\linewidth,height=2.5cm,multiline=true,
           bordercolor=black]{}

\vs

\noindent\textbf{D5.} Why does \texttt{git push} require a prior
\texttt{git commit}? What is the conceptual difference between
the two operations?
(See Chapter~\ref{ch:git}.)\\[4pt]
\TextField[name=d5,width=\linewidth,height=2.5cm,multiline=true,
           bordercolor=black]{}

\vs

\noindent\textbf{D6.} In \LaTeX{}, what is the practical difference
between \verb|\[...\]| and \verb|$$...$$|? Which should be used in
new documents, and why?
(See Chapter~\ref{ch:latex}, Section~\ref{sec:latex-math}.)\\[4pt]
\TextField[name=d6,width=\linewidth,height=2.5cm,multiline=true,
           bordercolor=black]{}

%================================================================
\paragraph{Optional: LLM Infrastructure (Ollama)}
\label{sec:assess-optional}
%================================================================

These items apply only to students who install and configure the
local Ollama LLM infrastructure. They are not required for the data
challenge.

\vs

\noindent\CheckBox[name=o1,width=0.4cm,height=0.4cm]{}~\textbf{O1.}
Ollama is installed and \texttt{ollama list} returns at least one
model in the terminal.
(See Chapter~\ref{ch:ollama-infrastructure},
Section~\ref{sec:install}.)

\vs

\noindent\CheckBox[name=o2,width=0.4cm,height=0.4cm]{}~\textbf{O2.}
\texttt{connectivity\_test.py} passes all four checks (TCP
handshake, model list, inference).
(See Chapter~\ref{ch:ollama-infrastructure}.)

\vs

\noindent\textbf{O3.} Explain the tradeoff between context window
size and memory consumption in Ollama. Why is \texttt{OLLAMA\_NUM\_CTX=8192}
a safe default for large models on hardware with limited memory?
(See Chapter~\ref{ch:ollama-infrastructure},
Section~\ref{sec:memory}.)\\[4pt]
\TextField[name=o3,width=\linewidth,height=3cm,multiline=true,
           bordercolor=black]{}

\end{Form}
